\subsubsection{Orthogonal Polynomials}\label{subsec:orthogonal-polynomials}

According to the Weierstrass approximation theorem, polynomials are dense in $L^2(-1,1)$. It is therefore natural to apply the Schmidt orthogonalization procedure to the sequence $1,x,x^2,\ldots$ and obtain a basis consisting of polynomials. We claim that up to factors these orthogonal polynomials are given by
\begin{equation}\label{eq:legendre-polynomial-rodrigues-formula}
P_n(x)=\dfrac1{2^nn!}\dfrac{\mathrm{d}^n(x^2-1)^n}{\mathrm{d}x^n}.
\end{equation}


First of all, it is obvious that $P_n$ is a polynomial of degree $n$. Moreover, integration by parts shows that
\begin{equation}\label{eq:legendre-lower-monomial-orthogonality}
\int_{-1}^1 P_n(x)x^m\,\mathrm{d}x=0
\end{equation}
for every $m<n$; hence we also have
\begin{equation}\label{eq:legendre-polynomial-pairwise-orthogonality}
\int_{-1}^1 P_n(x)P_m(x)\,\mathrm{d}x=0
\end{equation}
for $m<n$, i.e., the $P_n$ are orthogonal in $L^2(-1,1)$. The $P_n$ are called Legendre polynomials. The first few of them are
\begin{equation}\label{eq:first-five-legendre-polynomials}
\begin{gathered}
P_0(x)=1,\quad P_1(x)=x,\quad P_2(x)=\dfrac32x^2-\dfrac12,\\
P_3(x)=\dfrac52x^3-\dfrac32x,\quad P_4(x)=\dfrac{35}8x^4-\dfrac{15}4x^2+\dfrac38.
\end{gathered}
\end{equation}


The Legendre polynomials are not normalized. Using repeated integration by parts, one finds
\begin{equation}\label{eq:legendre-square-integral-reduction}
\int_{-1}^1 P_n^2(x)\,\mathrm{d}x=\dfrac{(2n)!}{2^{2n}(n!)^2}\int_{-1}^1(1-x^2)^n\,\mathrm{d}x.
\end{equation}
The integral of $(1-x^2)^n$ can be evaluated by observing that $(1-x^2)^n=(1-x)^n(1+x)^n$ and using repeated integration by parts. The final result is
\begin{equation}\label{eq:legendre-square-integral-value}
\int_{-1}^1 P_n^2(x)\,\mathrm{d}x=\dfrac2{2n+1}.
\end{equation}


A variety of other orthogonal polynomials are also important for applications. These polynomials are orthogonal in weighted $L^2$-spaces.

\begin{Definition}
\label{def:weighted-lp-two-space}
Let $\Omega$ be an open set in $\mathbb{R}^m$, and let $w$ be a continuous function from $\Omega$ to $\mathbb{R}^+$. Then we define
\begin{equation}\label{eq:weighted-continuous-lp-two-space}
\widetilde L_w^2(\Omega):=\left\{u\in C(\overline\Omega)\mid\int_\Omega w(\bm{x})|u(\bm{x})|^2\,\mathrm{d}\bm{x}<\infty\right\}.
\end{equation}
The inner product is defined by
\begin{equation}\label{eq:weighted-lp-two-inner-product}
(u,v)=\int_\Omega w(\bm{x})\overline{u(\bm{x})}v(\bm{x})\,\mathrm{d}\bm{x}.
\end{equation}
The space $L_w^2(\Omega)$ is the completion of $\widetilde L_w^2(\Omega)$.
\end{Definition}

For any weight function $w$ and any interval $(a,b)$, we can now define orthogonal polynomials by orthogonalizing the sequence $1,x,x^2,\ldots$ (provided of course, that $w$ is such that polynomials are in $L_w^2(\Omega)$). The following cases are particularly important:
